% 建议负责人：负责需求分析与应用价值的成员
\section{项目背景}

\subsection{问题描述与实际需求}

课程评价是高校了解教学效果的重要信息来源。
学生通常会在评价中同时谈到教师讲解、课程内容、作业量、考试范围、实验环境和学习收获。
例如“老师讲得很清楚，但是作业太多”既包含正面态度，也包含明确问题。
只统计评分或采用简单正负词匹配，容易丢失评价中的转折关系和具体原因。

教育文本分析相关研究也指出，学生开放性反馈适合从整体情感、句子情感、实体和方面等多个层次分析；
方面级情感分析在教育场景中的价值，正是帮助教师区分学生对课程、教师、作业或考核等不同对象的态度\citep{shaik2023educational,hajrizi2020aspect}。
基于这一思路，本项目不仅输出积极、中性、消极标签，还保留课程维度和原文证据。

现有课程评价处理主要存在以下困难：

\begin{itemize}
  \item \textbf{文本数量大。}期末集中收集评价时，教师和课程负责人难以逐条阅读并归纳。
  \item \textbf{表达形式复杂。}评价可能是中文、英文或中英混合文本，还可能包含缩写、拼写错误、
        否定、建议和讽刺等非规范表达。
  \item \textbf{情感与问题维度并不等价。}同一句话可以整体中性，但在作业任务或考试安排上包含
        明显负面反馈，仅输出情感标签不足以支持教学改进。
  \item \textbf{结果需要可解释。}教学场景不能只给出“负面”结论，还应说明命中了哪个课程维度、
        对应原文证据是什么，以及应该优先采取什么措施。
\end{itemize}

本项目希望将非结构化课程评价转换为结构化信息，包括语言、情感标签、置信度、课程维度、
关键词、证据片段、相似评价和改进建议，从而降低人工整理成本，为教学改进提供更直接的依据。

\subsection{使用场景与目标用户}

系统主要面向三类使用者：

\begin{enumerate}
  \item \textbf{任课教师：}分析单条典型评价，了解学生认可点和具体困难，辅助调整教学节奏、
        作业量和实验安排等。
  \item \textbf{课程负责人：}批量导入课程评价 CSV，查看情感分布、课程维度分布和高频关键词，
        识别需要优先处理的共性问题。
  \item \textbf{教学管理人员：}通过统一的结构化指标比较不同课程或不同学期的反馈趋势，并保留
        可追溯的原文证据。
\end{enumerate}

典型使用流程为：教师上传评价数据，系统进行双语预处理和情感分类，再按照教学内容、授课
方式、作业任务、考试安排、实验实践、学习收获六个维度归纳问题，最后生成便于教师阅读的总结
与建议。

\subsection{项目目标与边界}

本项目不替代教师做教学决策，而是构建一个可运行、可解释、可降级的 NLP 辅助
系统。具体目标如下：

\begin{itemize}
  \item 支持中文、英文和中英混合课程评价；
  \item 完成正面、中性、负面三分类，并处理常见混合评价；
  \item 将评价映射到课程维度，同时展示命中关键词和原文证据；
  \item 支持首页概览、单条分析、批量分析、固定案例验证和模型评估展示；
  \item 在 BERT、传统模型或外部 API 不可用时保持主要流程可运行；
  \item 提供可复现的数据处理、模型训练、消融实验和压力测试脚本。
\end{itemize}

项目在设计阶段明确区分三类结论：独立测试集用于比较模型性能，固定案例用于验证流程，独立压力测试用于观察复杂表达下的鲁棒性。
三类结果分别呈现，避免将某一组案例上的满分直接解释为系统整体准确率。

\noindent\textbf{项目成果概述}

本项目最终实现了一个可运行的课程评价分析系统，支持中文、英文和中英混合输入的单条评价
分析、批量 CSV 分析、固定案例验证和模型评估。系统综合使用 multilingual BERT、TF-IDF
Logistic Regression、规则校正和 LLM 建议生成，将课程评价转换为情感标签、课程维度、关键词、
原文证据、相似评价和教学建议。multilingual BERT 在 2280 条独立测试样本上的 Accuracy 为
0.7461，Macro-F1 为 0.7457；完整混合流程在 12 条固定案例上通过 12/12，在 48 条独立压力
测试上通过 34/48。项目同时设计了模型回退和本地模板兜底机制，以保证部署环境不稳定
时仍能运行。
